\section*{ID9895: Relation: P/kissingNumber\_max}
\paragraph{Metadata.}
PiID: 8700126456001 | RTEID: RTE:P38424.4825:T0.4799 | UUID: 07f7a320-64ff-58e5-ac4f-f2dcc2a54c87

\paragraph{Canonical Statement.}
log\_2(153600)\textbackslash{}approx17.2\$ bits and \$H\_\{\textbackslash{}text\{binary\}\}(\textbackslash{}pi - 3)\$ is some entropy of the fractional part of π (a constant \textasciitilde{}\$0.97\$ bits). So roughly \$C \textbackslash{}approx 18\$ qubits per use, indicating our channel can teleport states up to 18 qubits in size per operation (this is extremely high; typical quantum teleports are 1 qubit per entangled pair; our entangled resource is enormously larger, enabling batch teleportation). The channel’s noise model is special: we design Kraus operators \$K\_\{k\}\$ that correspond to errors where a state hops to a different kissing sphere of the same index (with probability \$p/kissingNumber\_\{max\}\$) or stays put (with probability \$1-p\$). This is like a tailored depolarizing channel that respects the kissing structure. Because of the stabilizer and topological codes (next step), we can correct those specific errors up to a threshold \$p < 1 - 1/\textbackslash{}pi \textbackslash{}approx 0.682\$. Error Correction (𝒞): On the receiving end, we apply a topological error-correcting code native to the kissing graph. In our implementation, we use a surface-code-like approach on a kissing lattice where stabilizers are defined by

\paragraph{Canonical Equation.}
\[
p/kissingNumber_{max}
\]

\paragraph{Dictionary.}
Relation: P/kissingNumber\_max — p/kissingNumber\_max

\paragraph{Encyclopedia.}
Relation: P/kissingNumber\_max is a symbolic expression, written p/kissingNumber\_max. It introduces a symbol or relation used elsewhere in the framework. It is built from p, k, s, n, g, N. Within the Trawin Zero Point registry it serves as one element of the framework's mathematical narrative.
